% Part: many-valued-logic
% Chapter: three-valued-logics
% Section: lukasiewicz

\documentclass[../../../include/open-logic-section]{subfiles}

\begin{document}

\olfileid{mvl}{inf}{god}

\olsection{د ګوډل منطقونه}

\begin{editorial}
  دا برخه د ګوډل ناطق ماتريس، صدق تابعې او له متناهي زنځيرونو او شهودي منطق سره ئې اړيکه بيانوي۔
\end{editorial}

\begin{defn}\ollabel{def:goedel} نامتناهي ارزښته د ګوډل
منطق~$\LogGod[\infty]$ په دې ماتريس تعريفېږي:
\begin{enumerate}
  \item معياري بياني ژبه~$\Lang L_0$ چې
  $\lfalse$, $\lnot$, $\land$, $\lor$, $\lif$ لري۔
  \item د رښتياوالي د قيمتونو سټ~$V_\infty$۔
  \item يوازې~$1$ ټاکل شوے قيمت دے؛ يعنې~$V^+ = \{1\}$۔
  \item صدق تابعې په لاندې تابعو ټاکل کېږي:
  \begin{align*}
    \tf{\lfalse} & = 0\\
    \tf{\lnot}[\LogGod](x) & = \begin{cases}
      1 & \text{که } x =0\\
      0 & \text{په بل حالت کښې}
    \end{cases}\\
    \tf{\land}[\LogGod](x,y) & = \min(x,y)\\
    \tf{\lor}[\LogGod](x,y) & = \max(x,y)\\
    \tf{\lif}[\LogGod](x,y) & = \begin{cases}
      1 & \text{که } x \le y\\
      y & \text{په بل حالت کښې.}
    \end{cases}
    \end{align*}
\end{enumerate}
د ګوډل $m$-ارزښته منطق همداسې تعريفېږي، خو~$V = V_m$ لري۔
\end{defn}

\begin{prop}
  د~\olref[thr][god]{defn:goedel} له مخې تعريف شوے منطق~$\LogGod[3]$
  د~\olref{def:goedel} له مخې له تعريف شوي~$\LogGod[3]$ سره يو شان دے۔
\end{prop}

\begin{proof}
  دا د~\olref[thr][god]{defn:goedel} د نښلوونکو صدق جدولونه
  د~\olref{def:goedel} له معادلو څخه له ترلاسه شوو جدولونو سره
  په پرتله کولو ليدل کېدے شي:
  \begin{center}
    \begin{tabular}{c|c} 
      $\tf{\lnot}[\LogGod[3]]$ & \\ 
      \hline  
      $1$ & $0$ \\ 
      $1/2$ & $0$ \\
      $0$ & $1$ 
    \end{tabular}
    \quad
    \begin{tabular}{c|ccc} 
      $\tf{\land}[\LogGod]$ & $1$ & $1/2$ & $0$ \\ 
      \hline 
      $1$ & $1$ & $1/2$ & $0$ \\ 
      $1/2$ & $1/2$ & $1/2$ & $0$\\ 
      $0$ & $0$ & $0$ & $0$ 
    \end{tabular}
    \\[2ex]
    \begin{tabular}{c|ccc} 
      $\tf{\lor}[\LogGod]$ & $1$ & $1/2$ & $0$ \\ 
      \hline 
      $1$ & $1$ & $1$ & $1$ \\ 
      $1/2$ & $1$ & $1/2$ & $1/2$ \\
      $0$ & $1$ & $1/2$ & $0$ 
    \end{tabular}
    \quad
    \begin{tabular}{c|ccc} 
      $\tf{\lif}[\LogGod]$ & $1$ & $1/2$ & $0$ \\ 
      \hline 
      $1$ & $1$ & $1/2$ & $0$ \\ 
      $1/2$ & $1$ & $1$ & $0$  \\ 
      $0$ & $1$ & $1$ & $1$ 
    \end{tabular}
  \end{center} 
\end{proof}

\begin{prop}\ollabel{prop:god-infty-m}
  که~$\Gamma \Entails[\LogGod[\infty]] !B$ وي، نو د هر~$m \ge 2$
  دپاره~$\Gamma \Entails[\LogGod[m]] !B$ هم وي۔
\end{prop}

\begin{proof}
  تمرين۔
\end{proof}

\begin{prob}
  \olref[mvl][inf][god]{prop:god-infty-m} ثابت کړئ۔
\end{prob}

که په~$\Gamma$ او~$!B$ کښې يوازې متناهي ډېر بياني متغيرونه
راشي، نو برعکس لورے هم صدق کوي۔ د يو نامتناهي مخالف مثال د
متغيرونو قيمتونه له صفر او يو سره يوځاے کړئ او په يوه متناهي
زنځير کښې په هماغه ترتيب کښېږدئ۔ د صفر، يو او ترتيب په ساتلو
سره نفي، عطف، فصل او شرط هم ساتل کېږي؛ نو مخالف مثال متناهي
ماتريس ته لېږدېږي۔
د نامتناهي مقدماتو لپاره عمومي برعکس لورے ناسم دے۔ د هر
$n\ge 0$ لپاره دا مقدمه واخلئ:
\[
(p_n \lif p_{n+1}) \lif q.
\]
که~$\pAssign v(q)<1$ وي او ټولې مقدمې پوره وي، د هر~$n$
لپاره~$\pAssign v(p_n)>\pAssign v(p_{n+1})$ پکار دي؛ که داسې
نۀ وي، دنننے شرط يو کېږي او بهرنۍ مقدمه د~$q$ هماغه قيمت
اخلي۔ په متناهي زنځير کښې دا نامتناهي نزولي لړۍ ناشونې ده۔
خو په ناطق ماتريس کښې~$\pAssign v(p_n)=1/(n+2)$ او
$\pAssign v(q)=1/2$ ټولې مقدمې پوره کوي، ځکه دننني شرطونه
د~$q$ تر قيمت زيات نۀ دي۔ دلته~$q$ پوره نۀ دے۔

د~$\LogGod[3]$ په شان، په~$\LogGod[\infty]$ کښې په شهودي منطق کښې
ټول معتبر !!{formula}s تاوتولوژۍ دي، او هماغه د ناتاوتولوژيو بېلګې
په~$\LogGod[\infty]$ کښې هم ناتاوتولوژۍ دي:
\begin{align*}
  & p \lor \lnot p && (p \lif q) \lif (\lnot p \lor q) \\
  & \lnot\lnot p \lif p && \lnot(\lnot p \land \lnot q) \lif (p \lor q) \\
  & ((p \lif q) \lif p) \lif p && \lnot(p \lif q) \lif (p \land \lnot q)
\end{align*}
يو داسې !!{formula} چې په شهودي منطق کښې معتبر نۀ دے، خو
په~$\LogGod[3]$ کښې تاوتولوژي ده، يعنې~$(p \lif q) \lor (q \lif
p)$، په~$\LogGod[\infty]$ کښې هم تاوتولوژي ده۔ په حقيقت کښې
$\LogGod[\infty]$ د شهودي منطق په توګه هم پېژندل کېدے شي چې
شېما~$(!A \lif !B) \lor (!B \lif !A)$ ورزياته شوې وي۔ دا مايکل
ډمېټ وښودله؛ نو~$\LogGod[\infty]$ ته ډېر ځله د ګوډل--ډمېټ
منطق~$\Log{LC}$ ويل کېږي۔

\begin{prob}
  وښيئ چې~$(p \lif q) \lor (q \lif p)$ د~$\LogGod[\infty]$
  تاوتولوژي ده۔
\end{prob}

\begin{prob}
  وښيئ چې~$(p \lif q) \lor (q \lif r) \lor (r \lif s)$، چې
  په~$\LogGod[3]$ کښې تاوتولوژي ده، په~$\LogGod[\infty]$
  کښې تاوتولوژي نۀ ده۔
\end{prob}

\end{document}
